% 习题库——第 4 章
% 章节由本文件的文件名确定。

\begin{exo}[沿路径积分微分形式]{integrale-forme-differentielle}
\keywords{微分形式, 恰当微分, 施瓦茨判据, 曲线积分}

考虑两个微分形式
\[
\omega_1 = y\,\mathrm{d}x + x\,\mathrm{d}y,
\qquad
\omega_2 = y\,\mathrm{d}x,
\]
以及从原点 $O(0,0)$ 连接到点 $M(1,1)$ 的三条路径：路径
$\gamma_1$ 先沿水平线段 $O\to(1,0)$，再沿竖直线段 $(1,0)\to M$；
路径 $\gamma_2$ 先沿竖直线段 $O\to(0,1)$，再沿水平线段 $(0,1)\to M$；
路径 $\gamma_3$ 则是直接连接 $O$ 与 $M$ 的直线段。

\begin{enumerate}
\item 通过找出满足 $\omega_1 = \mathrm{d}f$ 的函数 $f$，证明 $\omega_1$ 是恰当微分形式。
\item 分别参数化各线段，计算 $\int_{\gamma_1}\omega_1$ 和 $\int_{\gamma_2}\omega_1$。再不经计算得到该结果。
\item 对 $\omega_2$ 应用施瓦茨判据。可以得出什么结论？
\item 计算 $\int_{\gamma_1}\omega_2$、$\int_{\gamma_2}\omega_2$ 和 $\int_{\gamma_3}\omega_2$，并作评论。
\end{enumerate}

\begin{indication}
在水平线段上 $\mathrm{d}y = 0$，在竖直线段上 $\mathrm{d}x = 0$：
每个积分都化为一个普通定积分。
\end{indication}

\begin{solution}
\textbf{问题 1。} $f(x,y) = xy$ 符合要求，因为 $\partial f/\partial x = y$ 且 $\partial f/\partial y = x$。

\textbf{问题 2。} 对于参数化为 $t\mapsto\bigl(x(t),y(t)\bigr)$ 的曲线，有
\[
\int_\gamma\omega_1
= \int \left[y(t)x'(t)+x(t)y'(t)\right],\mathrm{d}t.
\]

路径 $\gamma_1$ 由两条线段组成。在第一段上，$x(t)=t$、$y(t)=0$，
$t\in[0,1]$；因此 $\mathrm{d}x=\mathrm{d}t$ 且 $\mathrm{d}y=0$。
在第二段上，$x(t)=1$、$y(t)=t$，同样有 $t\in[0,1]$；因此
$\mathrm{d}x=0$ 且 $\mathrm{d}y=\mathrm{d}t$。所以
\[
\begin{aligned}
I_1
= \int_{\gamma_1}\omega_1
&= \int_0^1\left(0\times 1+t\times 0\right)\,\mathrm{d}t
 + \int_0^1\left(t\times 0+1\times 1\right)\,\mathrm{d}t \\
&= 0+\int_0^1\mathrm{d}t = 1.
\end{aligned}
\]

对于 $\gamma_2$，第一段参数化为 $x(t)=0$、$y(t)=t$，第二段参数化为
$x(t)=t$、$y(t)=1$，两段均有 $t\in[0,1]$。因此
\[
\begin{aligned}
I_2
= \int_{\gamma_2}\omega_1
&= \int_0^1\left(t\times 0+0\times 1\right)\,\mathrm{d}t
 + \int_0^1\left(1\times 1+t\times 0\right)\,\mathrm{d}t \\
&= 0+\int_0^1\mathrm{d}t = 1.
\end{aligned}
\]
两个数值相同。其实不计算也可以预见这一点：由于 $\omega_1=\mathrm{d}f$ 且
$f(x,y)=xy$，其积分只依赖于端点，
\[
\int_\gamma\omega_1=f(M)-f(O)=f(1,1)-f(0,0)=1.
\]

\textbf{问题 3。} 写成 $\omega_2=A(x,y)\,\mathrm{d}x+B(x,y)\,\mathrm{d}y$。
这里 $A(x,y)=y$、$B(x,y)=0$。若该形式恰当，施瓦茨判据应给出
\[
\frac{\partial A}{\partial y}=\frac{\partial B}{\partial x}.
\]
然而 $\partial A/\partial y=1$，而 $\partial B/\partial x=0$。
两个交叉偏导数不同，因此 $\omega_2$ 不是恰当微分形式。它的积分可能依赖于
$O$ 与 $M$ 之间所取的路径，下面的计算正好验证了这一点。

\textbf{问题 4。} 因为 $\omega_2=y\,\mathrm{d}x$，只有在 $y$ 非零时
$x$ 的变化才会对积分有贡献。沿用前面的参数化，在 $\gamma_1$ 上有
\[
\begin{aligned}
J_1
=\int_{\gamma_1}\omega_2
&=\int_0^1 0\times 1\,\mathrm{d}t
  +\int_0^1 t\times 0\,\mathrm{d}t \\
&=0.
\end{aligned}
\]
在 $\gamma_2$ 上，第一段满足 $x(t)=0$，故 $\mathrm{d}x=0$；
第二段满足 $x(t)=t$、$y(t)=1$ 且 $\mathrm{d}x=\mathrm{d}t$。于是
\[
\begin{aligned}
J_2
=\int_{\gamma_2}\omega_2
&=\int_0^1 t\times 0\,\mathrm{d}t
  +\int_0^1 1\times 1\,\mathrm{d}t \\
&=1.
\end{aligned}
\]
最后，对角线段 $\gamma_3$ 可显式参数化为
\[
x(t)=t,\qquad y(t)=t,\qquad t\in[0,1],
\]
从而 $\mathrm{d}x=\mathrm{d}t$、$\mathrm{d}y=\mathrm{d}t$。故
\[
J_3
=\int_{\gamma_3}\omega_2
=\int_0^1 y(t)x'(t)\,\mathrm{d}t
=\int_0^1 t\,\mathrm{d}t
=\frac12.
\]
三条路径的端点相同，却给出三个不同的值：$J_1=0$、$J_2=1$ 和
$J_3=1/2$。这正是非恰当微分形式所特有的路径依赖性。
\end{solution}
\end{exo}

\begin{exo}[相同的内能变化，三种传递方式]{meme-delta-u-trois-transferts}
\keywords{热力学第一定律, 内能, 热, 机械功, 电功, 量热学, 状态函数}

一种液体、其量热器以及一个浸没其中的小电阻组成一个宏观静止的封闭系统。
容器是刚性的，并假设系统的总热容恒定：
\[
C=2{,}00\ \mathrm{kJ\,K^{-1}}.
\]
要使系统从温度为 $T_A=20{,}0\,^\circ\mathrm{C}$ 的平衡态 $A$
变到温度为 $T_B=25{,}0\,^\circ\mathrm{C}$ 的平衡态 $B$。已知
\[
U(B)-U(A)=C(T_B-T_A).
\]
宏观动能和宏观势能的变化均为零。忽略向外界的一切损失。

从同一初态 $A$ 出发，彼此独立地实施以下三个方案：
\begin{itemize}
\item[方案 1] 使系统与一个高温热源接触，但不向系统提供任何功。
\item[方案 2] 容器壁绝热。一个质量为 $m$ 的重物从高度 $h=5{,}00$ m 落下，带动桨叶搅拌液体。最后桨叶和液体均静止，重物损失的全部机械能都传递给系统。取 $g=9{,}81\ \mathrm{m\,s^{-2}}$。
\item[方案 3] 系统获得热量 $Q_3=4{,}00$ kJ，其余能量通过电路供给电阻。系统获得的电功率恒定，为 $\mathcal P=300$ W。
\end{itemize}

\begin{enumerate}
\item 计算三个方案共同的内能变化 $\Delta U=U(B)-U(A)$。
\item 对前两个方案，分别确定系统获得的热量 $Q$ 和功 $W$。在第二个方案中计算所需的质量 $m$。
\item 对第三个方案，计算系统获得的电功以及电阻通电的时间。
\item 三个过程的初态和终态相同。解释为什么它们的 $Q$ 和 $W$ 仍可不同。系统在终态 $B$ 中是否含有“热”或“功”？
\end{enumerate}

\begin{indication}
采用银行家符号约定，应用 $\Delta U=Q+W$。对于桨叶，系统获得的功等于
重物势能的减少量 $mgh$。通过导线穿过系统边界的电能计作电功。
\end{indication}

\begin{solution}
\textbf{问题 1。} 温差为 $T_B-T_A=5{,}0$ K。因此
\[
\Delta U=C(T_B-T_A)
=2{,}00\times5{,}0
=10{,}0\ \mathrm{kJ}.
\]
该数值只取决于两个平衡态 $A$ 和 $B$。

\textbf{问题 2。} 在方案 1 中，系统没有获得功：$W_1=0$。
由第一定律得
\[
Q_1=\Delta U=10{,}0\ \mathrm{kJ}.
\]
热量为正，因为能量进入系统。

在方案 2 中，容器壁绝热，故 $Q_2=0$。全部内能变化都来自桨叶的机械功：
\[
W_2=\Delta U=10{,}0\ \mathrm{kJ}.
\]
又因 $W_2=mgh$，所以
\[
m=\frac{W_2}{gh}
=\frac{10{,}0\times10^3}{9{,}81\times5{,}00}
\simeq 204\ \mathrm{kg}.
\]
这个较大的数量级提醒我们，即使温度仅小幅升高，也已经对应相当大的机械能。

\textbf{问题 3。} 第一定律要求
\[
W_3=\Delta U-Q_3=10{,}0-4{,}00=6{,}00\ \mathrm{kJ}.
\]
该功为正：电源向系统提供能量。利用 $W_3=\mathcal P\,\Delta t$，有
\[
\Delta t=\frac{W_3}{\mathcal P}
=\frac{6{,}00\times10^3}{300}
=20{,}0\ \mathrm{s}.
\]

\textbf{问题 4。} 三条路径分别给出
\[
(Q_1,W_1)=(10{,}0,0)\ \mathrm{kJ},
\qquad
(Q_2,W_2)=(0,10{,}0)\ \mathrm{kJ},
\]
以及
\[
(Q_3,W_3)=(4{,}00,6{,}00)\ \mathrm{kJ}.
\]
在所有情况下，$Q+W$ 均等于 $10{,}0$ kJ，并产生相同的变化 $\Delta U$。
内能是状态函数，而热和功描述的是过程中发生的能量传递。在终态 $B$ 中，
系统具有内能，但并不含有“热”或“功”。
\end{solution}
\end{exo}

\begin{exo}[恒定外压下的压缩]{compression-pression-exterieure-constante}
\seoready{true}
\keywords{压强力的功, 恒定外压, 压缩, 膨胀, 真空膨胀}

气体在恒定外压 $P_0$ 下从体积 $V_A$ 压缩到体积 $V_B<V_A$。

\begin{enumerate}
\item 计算气体获得的功。
\item 气体克服同一外压 $P_0$，从 $V_B$ 回到 $V_A$。计算气体获得的功，并与压缩过程的功比较。
\item 当气体在真空中从 $V_B$ 膨胀到 $V_A$ 时，重新讨论该膨胀过程。解释所得的三个符号。
\end{enumerate}

\begin{indication}
使用 $W=-\int P_{\mathrm{ext}}\,\mathrm{d}V$。必须对\emph{外部}压强积分，
即使过程不是准静态过程也如此。
\end{indication}

\begin{solution}
\textbf{问题 1。} 外压恒定，因此
\[
W_{A\to B}=-P_0(V_B-V_A)=P_0(V_A-V_B)>0.
\]
压缩时气体获得功。

\textbf{问题 2。} 对返回过程，
\[
W_{B\to A}=-P_0(V_A-V_B)<0.
\]
两个功互为相反数，因为两个过程都是克服同一恒定外压进行的。

\textbf{问题 3。} 在真空中，$P_{\mathrm{ext}}=0$，因此
\[
W_{B\to A}^{\mathrm{vide}}=0.
\]
所以膨胀并不自动意味着功为负：必须确实有外力被气体推开。
\end{solution}
\end{exo}
